% GENERATED FILE. Edit canonical lessons, then run npm run generate:books.
% canonical-source-hash: fnv1a64:53fb380110e14aa1
% canonical-lessons: PA-W10-a1-con-reloj

\chapter{Writing --- With A Clock}
\label{ch:pa-ghadi-naal}
\hlchaptermodality{46}

\begin{chapteropening}
\textbf{By the end of this chapter:} \emph{I can complete the A1 writing paper's two tasks inside its published twenty minutes, filling every form field and keeping the message between thirty and forty words.}

Complete PA-W10-a1-con-reloj: seven minutes for the form and thirteen for the message, then count the words and check the form for blanks before mistakes.
\end{chapteropening}

\section[(timed)]{See the A1 writing clock}

\label{lesson:PA-W10-a1-con-reloj}



\begin{quote}
You have filled some familiar fields with as long as you liked.

Read the later paper's conditions now. This short lesson only explains them; it does not attempt or score either task under a clock.
\end{quote}

\subsection*{Writing — the published paper envelope}
\textbf{What the contract publishes.} This track's A1 writing paper runs \textbf{20 minutes} and holds \textbf{two} tasks:

\begin{itemize}
\raggedright
  \item a practical form asking for \textbf{six to eight} completed fields, worth 40 of the paper's 100 points;
  \item a message for a named reader and a stated purpose, \textbf{30 to 40 words}, worth the other 60.
\end{itemize}

Those numbers are read out of \texttt{task-shapes/a1.json}, the same file the assessment contract points at. They are not invented here.

Notice what the split says before you write a word: the \textbf{message is worth more than the form}, and it is the shorter of the two. Time spent perfecting a field is time taken from the part that scores.

One possible whole-paper plan gives the form about seven minutes, but that is a planning example, not a timer starting in this lesson. First, learn every field and practise all six together without a model.

The phone field is the one to watch under a clock. You have practised grouping its digits; grouping is what makes a number readable back, and it is the first thing a hurried hand drops.

The remaining thirteen minutes would be for the 30–40-word message to one named reader. That message also needs its own untimed runway first.

Later short timed checkpoints will train each task separately. Two complete mocks will test the full twenty-minute paper after that.

\begin{tcolorbox}[breakable,colback=teal!4,colframe=teal!35!black,title={Before you move on}]
Recall the message's word band: \textbf{30–40 words}. You are not writing that message yet.

Recall the form's size: \textbf{six to eight fields}. Those fields must be taught before they can be scored.

This orientation awards no timed-writing credit. Timed production comes only after independent form and named-reader message writing are both secure without a visible answer model.
\end{tcolorbox}
